%! Principles of Experimental Design — Unit 1, Lesson 1.6 — Slides (Beamer)
% PILOT SHAPE (2026-09-06): modeled on AP Statistics lesson 1.4 minus its AP Practice page.
% GRADUAL RELEASE deck: title → targets → warm-up → hook (unresolved) → I DO divider → two or
% three frames per notes ROW (labelled by row name and problem numbers) → WE DO (a live
% surface, un-answered) → spoken debrief with the cold check → close & START THE HOMEWORK,
% which is the individual practice. No group activity, no practice launch, no exit ticket.
% There is no spoiler rule: the notes frames ARE the direct instruction.
% saar-beamer does NOT load saar-article (so \CourseName is undefined — the course name is
% literal) and does NOT load tcolorbox (use block); beamer itemize takes no [leftmargin=…].
\documentclass[aspectratio=169, 11pt]{beamer}
\usepackage{saar-beamer}   % provides \ceruleanheader{} and \sectionlabel[color]{}

\newcommand{\redink}[1]{{\color{redacc}\bfseries #1}}

\begin{document}

% TITLE SLIDE (CourseName is not defined in beamer, so the name is literal)
{
\setbeamercolor{background canvas}{bg=cerulean}
\begin{frame}[plain]
  \vspace{2.8cm}\hspace{0.55cm}%
  \begin{minipage}{0.9\textwidth}
    {\color{goldacc}\bfseries\small LESSON 1.6}\\[0.5cm]
    {\color{white}\bfseries\LARGE Principles of Experimental Design}\\[1.2cm]
    {\color{frostmid}\small Statistical Analysis \& Algebraic Reasoning~$\cdot$~Unit 1}
  \end{minipage}
\end{frame}
}

% ── LEARNING TARGETS ─────────────────────────────────────────────────────────
\begin{frame}[t, plain]
  \ceruleanheader{Today's targets}
  \sectionlabel{Lesson 1.6}
  By the end of today, I can\ldots
  \begin{itemize}\setlength{\itemsep}{5pt}
    \item tell an \textbf{experiment} from an observational study by asking \emph{who built
          the groups} --- and name the \textbf{units}, the \textbf{treatment}, and the
          \textbf{control group}.
    \item check a design against \textbf{comparison, randomization, replication, control} ---
          and name the principle a changed plan breaks.
    \item explain the \textbf{placebo effect}, say who is \textbf{blinded}, and split units
          into \textbf{blocks}.
    \item justify why \textbf{random assignment} lets a study say \emph{caused} --- and why
          the size of the gap does not.
  \end{itemize}
  \begin{block}{How today works}
    Warm-up $\to$ \textbf{notes together, row by row} $\to$ \textbf{one experiment we work as
    a class} $\to$ debrief $\to$ \textbf{you start the homework here, alone}.
  \end{block}
\end{frame}

% ── WARM-UP ──────────────────────────────────────────────────────────────────
\begin{frame}[t, plain]
  \ceruleanheader{Warm-Up --- 5 minutes, silent}
  \sectionlabel{Riverbend High School --- this time they assign}
  \vspace{0.2cm}
  \begin{center}
    \begin{minipage}{0.88\textwidth}
      \textbf{120} students volunteered. A random number generator split them into two groups
      of \textbf{60}: one attended the Study Center twice a week, the other not at all.
      \begin{enumerate}\setlength{\itemsep}{6pt}
        \item \textbf{45} of the \textbf{60} who attended passed every class. What percent?
        \item \textbf{33} of the \textbf{60} who did not attend passed every class. What percent?
        \item[3(a).] How many \emph{percentage points} higher is the group that attended?
        \item[3(b).] Last quarter she only \emph{watched}: \textbf{65\%} of Study Center users
              passed, against \textbf{80\%} of those who never went. Which group was higher
              then, and by how much?
      \end{enumerate}
    \end{minipage}
  \end{center}
  \vspace{0.3cm}
  \begin{center}
    {\color{cerulean}\bfseries Both denominators are 60. That is not an accident --- keep the page.}
  \end{center}
\end{frame}

% ── HOOK — leave it UNRESOLVED; problem 12 (row 4, Caused) settles it ────────
{
\setbeamercolor{background canvas}{bg=cerulean}
\begin{frame}[plain]
  \vspace{1.1cm}
  \begin{center}
    {\color{frostmid}\bfseries\small SAME SCHOOL. SAME STUDY CENTER. TWO QUARTERS.}\\[0.7cm]
    {\color{white}\bfseries\Large Watched: users 65\% \quad vs.\quad non-users 80\%}\\[0.35cm]
    {\color{white}\bfseries\Large Assigned: treated 75\% \quad vs.\quad control 55\%}\\[0.8cm]
    {\color{goldacc}\bfseries\large Opposite conclusions. Which study should she believe?}\\[0.45cm]
    {\color{frostmid}\small Vote: last quarter's, this quarter's, or neither --- and write your
    reason. We settle it at problem 12.}
  \end{center}
\end{frame}
}

% ── I DO — direct instruction begins ─────────────────────────────────────────
{
\setbeamercolor{background canvas}{bg=cerulean}
\begin{frame}[plain]
  \vspace{1.8cm}\hspace{0.8cm}%
  \begin{minipage}{0.86\textwidth}
    {\color{goldacc}\bfseries\small GUIDED NOTES \ \textbullet\ \ I DO}\\[0.7cm]
    {\color{white}\bfseries\Large Open to your Guided Notes page. We work the table row by
    row --- fill each term as we name it, not before.}\\[0.9cm]
    {\color{frostmid}\small Five terms go in the vocabulary box today. Write each one the
    moment we use it.}
  \end{minipage}
\end{frame}
}

% ── ROW 1a. ASSIGN, DON'T WATCH ──────────────────────────────────────────────
\begin{frame}[t, plain]
  \ceruleanheader{Assign --- don't just watch}
  \sectionlabel{notes row 1 $\cdot$ Experiment $\cdot$ problems 1--3}
  \vspace{4pt}
  \begin{columns}[t]
    \begin{column}{0.47\textwidth}
      \begin{block}{Observational study (Lesson 1.5)}
        You \textbf{watch}. The individuals were already in their groups --- they chose.
      \end{block}
      \vspace{4pt}
      {\small Riverbend, last quarter: 150 records, students chose the Study Center.}
    \end{column}
    \begin{column}{0.47\textwidth}
      \begin{block}{Experiment}
        You \textbf{assign} the units to groups and \textbf{impose a treatment} on at least one.
      \end{block}
      \vspace{4pt}
      {\small Riverbend, this quarter: 120 volunteers, a generator, 60 and 60.}
    \end{column}
  \end{columns}
  \vspace{10pt}
  \begin{center}
    {\large The test is one question: \redink{who built the groups?}}\\[6pt]
    {\small\color{slate} The 120 volunteers are the \textbf{experimental units}
    (\textbf{subjects}). Twice-a-week attendance is the \textbf{treatment}. The other 60 are
    the \textbf{control group}.}
  \end{center}
\end{frame}

% ── ROW 1b. THE TRAP — problem 4 ─────────────────────────────────────────────
\begin{frame}[t, plain]
  \ceruleanheader{Problem 4 --- vote before you write}
  \sectionlabel{notes row 1 $\cdot$ the trap}
  \vspace{6pt}
  \begin{center}
    \begin{minipage}{0.86\textwidth}
      ``Give \textbf{60} randomly chosen volunteers a new tutoring app and the other
      \textbf{60} the old worksheet packet.''
      \vspace{6pt}
      \begin{center}\large Is there a \textbf{control group}?\end{center}
    \end{minipage}
  \end{center}
  \vspace{8pt}
  \begin{block}{The room usually says no --- everybody got something}
    The \textbf{control group is the baseline you compare against}, not the group that gets
    nothing. The packet group \redink{is} the control group. Take away the comparison and the
    first principle is gone.
  \end{block}
\end{frame}

% ── ROW 2. THE FOUR PRINCIPLES ───────────────────────────────────────────────
\begin{frame}[t, plain]
  \ceruleanheader{Four principles --- every experiment worth trusting}
  \sectionlabel{notes row 2 $\cdot$ Principles $\cdot$ problems 5--7}
  \vspace{2pt}
  \begin{center}
    \small
    \renewcommand{\arraystretch}{1.3}
    \begin{tabular}{@{} l l @{}}
    \hline
    \textbf{Principle} & \textbf{What it means} \\
    \hline
    Comparison     & At least two groups, one of them a control group. \\
    Randomization  & Chance decides the groups --- not the researcher, not the subject. \\
    Replication    & Many units in each group, so one person's luck cannot decide it. \\
    Control        & Every other condition the same for both groups. \\
    \hline
    \end{tabular}
  \end{center}
  \vspace{8pt}
  \begin{block}{Problem 7 --- Cedar Ridge}
    The reminder group also got \textbf{new recycling bins}. That is a
    \textbf{\color{cerulean}confounding variable}: it differs between the groups \emph{and}
    affects the response, so nobody can say which one worked.
    \redink{Randomization makes the groups alike before you start; control keeps them alike
    while you run.}
  \end{block}
\end{frame}

% ── ROW 3. BLINDING AND BLOCKING ─────────────────────────────────────────────
\begin{frame}[t, plain]
  \ceruleanheader{Two refinements: blinding and blocking}
  \sectionlabel{notes row 3 $\cdot$ Blinding \& Blocking $\cdot$ problems 8--9}
  \vspace{4pt}
  \begin{columns}[t]
    \begin{column}{0.47\textwidth}
      \begin{block}{Blinding}
        \textbf{Placebo effect:} a subject responds to the \emph{idea} of being treated.
        A \textbf{placebo} is a fake treatment built to look real.
        \textbf{Single-blind:} the subjects don't know. \textbf{Double-blind:} the subjects
        \emph{and} the people scoring don't know.
      \end{block}
      \vspace{3pt}
      {\small Room A tutoring, Room B quiet study --- \emph{both} meet after school.}
    \end{column}
    \begin{column}{0.47\textwidth}
      \begin{block}{Blocking}
        Sort the units into \textbf{blocks} of similar units \emph{first}, then randomize
        \textbf{inside} each block. \textbf{Matched pairs} is the block of two.
      \end{block}
      \vspace{3pt}
      {\small Block on last quarter's grades: 80 passed $\to$ 40 / 40; \quad 40 did not $\to$
      20 / 20; \quad \textbf{60 per group}.}
    \end{column}
  \end{columns}
  \vspace{10pt}
  \begin{center}
    {\small\color{slate} Riverbend's plain 60/60 split is a \textbf{completely randomized
    design}. Blocking does not replace chance --- it applies chance \emph{inside} groups of
    similar units.}
  \end{center}
\end{frame}

% ── ROW 4. THE CRUX ──────────────────────────────────────────────────────────
\begin{frame}[t, plain]
  \ceruleanheader{The design buys the verb --- not the gap}
  \sectionlabel[redacc]{notes row 4 $\cdot$ Caused $\cdot$ problems 10--13 $\cdot$ the whole point of today}
  \vspace{4pt}
  \begin{center}
    \small
    \renewcommand{\arraystretch}{1.3}
    \begin{tabular}{@{} l c c l @{}}
    \hline
    \textbf{Riverbend} & \textbf{Who chose?} & \textbf{Gap} & \textbf{May claim} \\
    \hline
    Last quarter --- watched  & the students & 15 points (the wrong way) & an association \\
    This quarter --- assigned & nobody: chance & 20 points & \textbf{caused} \\
    \hline
    \end{tabular}
  \end{center}
  \vspace{6pt}
  \begin{itemize}\setlength{\itemsep}{4pt}
    \item \textbf{Problem 12 --- vote again:} last quarter's, this quarter's, or neither?
          ``20 points beats 15'' --- is that the reason? \redink{No.}
    \item \textbf{Problem 13:} suppose the watched study had found 90\% vs.\ 55\% --- a
          \textbf{35}-point gap. Believe it now? \redink{Still no.}
  \end{itemize}
  \begin{block}{Why (PS.DC.3b)}
    Chance built the two groups, so they were alike \emph{before} the treatment --- the
    treatment is the only thing left that can explain the gap. A self-chosen gap is inflated
    by \emph{whoever chose}.
  \end{block}
\end{frame}

% ── WE DO — a LIVE surface: task and data only, no answers ───────────────────
\begin{frame}[t, plain]
  \ceruleanheader{Guided Practice: Harbor Point Swim Times}
  \sectionlabel[goldacc]{We do $\cdot$ problems 14--17 $\cdot$ you hold the pen}
  Last lesson the pool \emph{watched}: 70\% of morning swimmers slept 7+ hours against 45\% of
  evening swimmers --- but members chose their own times, and 75\% of the morning swimmers were
  retired. Now the director \emph{assigns}: \textbf{90} volunteers, a generator, \textbf{45} at
  a reserved 6 a.m.\ slot and \textbf{45} at 7 p.m., four weeks. \textbf{27} and \textbf{18}
  averaged 7+ hours.
  \begin{itemize}\setlength{\itemsep}{4pt}
    \item[14.] Units, treatment, control group --- and who decided the slots.
    \item[15.] The percent for each group, and how many points apart.
    \item[16.] The four principles, one phrase each. Then: who \emph{can} be blinded here?
    \item[17.] Watching gave 25 points; assigning gives 20. Which is the stronger evidence?
  \end{itemize}
  \begin{block}{The question behind the question}
    If you picked the 25 in problem 17, you are still at the hook vote.
  \end{block}
\end{frame}

% ── GUIDED PRACTICE problem 17 — shown AFTER the class has worked it ─────────
\begin{frame}[t, plain]
  \ceruleanheader{Why the smaller gap is the stronger evidence}
  \sectionlabel[goldacc]{Guided Practice, problem 17}
  \vspace{4pt}
  \begin{columns}[t]
    \begin{column}{0.47\textwidth}
      \begin{block}{Watching: 70\% vs.\ 45\% --- 25 points}
        \redink{Inflated.} 75\% of the morning swimmers were retired --- and retired members
        sleep more anyway. Part of those 25 points is retirement, not swim time.
      \end{block}
    \end{column}
    \begin{column}{0.47\textwidth}
      \begin{block}{Assigning: 60\% vs.\ 40\% --- 20 points}
        \redink{Earned.} Chance spread the retired members across both groups, so the 20
        points are the swim time by itself.
      \end{block}
    \end{column}
  \end{columns}
  \vspace{10pt}
  \begin{center}
    {\large A \redink{smaller} gap, honestly built, beats a bigger gap that somebody
    self-selected into.}\\[5pt]
    {\small\color{slate} Also true and worth saying: you cannot blind a swimmer to the time of
    day. The \emph{sleep-log scorers} can be blinded --- and that is the honest answer to
    problem 16.}
  \end{center}
\end{frame}

% ── DEBRIEF ──────────────────────────────────────────────────────────────────
\begin{frame}[t, plain]
  \ceruleanheader{Debrief: who built the groups?}
  \sectionlabel{10 minutes $\cdot$ pens down, papers up --- we say these aloud}
  \vspace{4pt}
  \begin{enumerate}\setlength{\itemsep}{4pt}
    \item The two Riverbend studies --- say the sentence. Why does the experiment earn
          \emph{caused}?
    \item Problem 4: the packet group got something. Is it still the control group?
    \item Problem 7: name the confounding variable at Cedar Ridge.
  \end{enumerate}
  \vspace{4pt}
  \begin{block}{Cold check --- hands down, thirty seconds}
    Bayside Middle School: 80 volunteers, a generator sends 40 to a daily reading block and 40
    to regular homeroom. Afterward 28 of the reading-block students improved and 16 of the
    homeroom students did. A teacher says: \emph{``Next year, just give the reading block to
    the students who want it.''} \textbf{Which principle does that break, and what could go
    wrong?}
  \end{block}
\end{frame}

% ── YOU DO — the homework is the individual practice, started in class ───────
\begin{frame}[t, plain]
  \ceruleanheader{You do: start the homework}
  \sectionlabel[goldacc]{Last 10 minutes $\cdot$ on your own $\cdot$ finish it at home}
  \begin{itemize}\setlength{\itemsep}{5pt}
    \item \textbf{Millbrook Public Library} --- 1{,}200 card holders; 200 volunteers, a
          generator, 100 in the Summer Reading Program and 100 on a fall waiting list.
    \item Name the parts; the two rates and the gap; pool the 200 and scale to 1{,}200
          \emph{as an estimate}.
    \item The same library, two summers --- one may claim an association, the other
          \emph{caused it}. Block on age and say why she blocked.
    \item Last summer's gap was 48 points; this summer's is 30. Explain why the
          \textbf{smaller} gap is the stronger evidence.
  \end{itemize}
  \begin{block}{Silent and alone --- I am circulating. Packet pages, scored out of 12, due at the start of the first class after two study halls.}
    Start with \textbf{1, 4, and 6}. Item 6 is the one I am reading over your shoulder --- if
    you picked the 48, flag me before you leave.
  \end{block}
\end{frame}

% ── CLOSE ────────────────────────────────────────────────────────────────────
{
\setbeamercolor{background canvas}{bg=cerulean}
\begin{frame}[plain]
  \vspace{1.5cm}\hspace{0.8cm}%
  \begin{minipage}{0.86\textwidth}
    {\color{goldacc}\bfseries\small BEFORE YOU GO}\\[0.7cm]
    {\color{white}\bfseries\Large An observational study watches. An experiment assigns --- and
    because chance built the groups, it is the one kind of study that earns the word
    \emph{caused}. The biggest gap is usually the one that earned it least.}\\[0.9cm]
    {\color{frostmid}\small Next: Lesson 1.7 --- Millbrook \emph{could} assign. Nobody can
    assign a person to smoke for thirty years. When is an experiment worth the cost, and when
    is it impossible or wrong to run one?}
  \end{minipage}
\end{frame}
}

\end{document}
